\section{Conclusion}
\label{sec:conclusion}

In this paper, we presented an end-to-end Video Foundation Model framework for fine-grained sports action recognition across 82 categories. By harnessing **VideoMAE ViT-B** with $16 \times 16 \times 2$ spatiotemporal tubelet embeddings and global self-attention across 12 transformer layers, our method captures non-linear athletic kinematics and resolves subtle fine-grained sports ambiguities. To eliminate experimental bias, we introduced a strict group-wise dataset partitioning protocol that eradicates scene, athlete, and camera leakage across 70\% train, 15\% validation, and 15\% test splits. Furthermore, our hardware-efficient cloud training architecture demonstrates that production-grade video transformers can be fully trained and deployed within free-tier cloud environments (Kaggle dual-T4 GPUs with Distributed Data Parallelism and Google Colab with persistent Google Drive checkpoints). Rigorous benchmarking against an extensive baseline hierarchy (Frozen Linear Probe, X3D-S, SlowFast-R50, and Video Swin-T) and thorough ablation studies demonstrate that VideoMAE establishes superior accuracy (91.8\% Top-1, 98.4\% Top-5) and notable sample efficiency. Future work will investigate cross-modal skeleton fusion (incorporating lightweight 3D pose heatmaps) and temporal action localization in continuous broadcast footage.
